%% AutomotiveUI 2026 — Work-in-Progress submission
%% ACM Primary Article Template, single-column (manuscript mode)
%% Anonymised for double-blind review

\documentclass[manuscript,screen,review,anonymous]{acmart}

%% Suppress ACM copyright / rights block for review copy
\setcopyright{none}
\acmConference[AutomotiveUI '26]{ACM International Conference on Automotive User Interfaces and Interactive Vehicular Applications}{September 20--23, 2026}{Gothenburg, Sweden}

%% ---------------------------------------------------------------
%% Packages
%% ---------------------------------------------------------------
\usepackage{booktabs}   % professional table rules
\usepackage{listings}   % code/YAML listings
\usepackage{xcolor}
\usepackage{graphicx}
\usepackage{hyperref}

\lstset{
  basicstyle=\small\ttfamily,
  breaklines=true,
  frame=single,
  rulecolor=\color{gray!40},
  backgroundcolor=\color{gray!5},
  xleftmargin=6pt,
  xrightmargin=6pt,
}

%% ---------------------------------------------------------------
%% Metadata
%% ---------------------------------------------------------------
\title{Toward a Semantic Mediation Layer for In-Vehicle Interaction}
\subtitle{An architecture of meaning that survives the displays, the agents, and the decade}

%% Authors suppressed for anonymous review
\author{Anonymous Author(s)}
\affiliation{\institution{Anonymous Institution}}
\email{anonymous@example.com}

\begin{abstract}
The software-defined vehicle is reshaping every layer of the in-cabin stack ---
except the one that matters most to the occupant. Industry effort concentrates
on pixels, frameworks, and ever-larger displays, while no shared abstraction
defines what an in-vehicle interaction \emph{means}, independent of which
screen renders it or which AI agent emits it. We argue that the SDV stack is
missing a \emph{mediation boundary}: a vendor-neutral layer in which
interactions are described by their meaning, attention demand, contextual
fitness, and authority of origin, rather than by buttons, screens, or widgets.
This paper sketches such a layer --- a Semantic Interaction Architecture
(SIA) --- with a typed node taxonomy (Actions, Events, States, Tasks),
declarative metadata contracts encoding predicted attention demand and trust
provenance, and a policy architecture for context-driven renderer selection.
We trace a safety-critical alert through verification, context-aware
translation, and adversarial rejection. The architecture opens concrete
handles for AutomotiveUI research: empirical calibration of attention
metrics, cross-modal equivalence studies, and trust calibration in
AI-augmented vehicles. It also identifies a near-term specification
agenda --- ontology formalism, multi-hop trust provenance, and
certified-vs-inferred acknowledgement --- that this community is positioned
to drive. The argument is provocative but structurally simple: existing
standards cover what lives \emph{below} the interaction (signals, transports,
services) and what lives \emph{above} it (renderers, widgets, frameworks).
The thing in the middle --- \emph{what an interaction is, in itself} --- has
no name, no contract, and no audit trail. SIA proposes one.
\end{abstract}

\keywords{software-defined vehicle, semantic interaction, HMI architecture,
  trust policy, attention model, multimodal rendering, in-vehicle agents}

\begin{document}
\maketitle

%% ---------------------------------------------------------------
\section{Introduction}
%% ---------------------------------------------------------------

The software-defined vehicle is reshaping every layer of the in-cabin stack
except the one that matters most to the occupant. Industry investment flows
into pixels, GPUs, larger displays, voice models, gesture recognition, and
AI agents --- each fighting at the surface. But there is no shared abstraction
that defines what an in-vehicle interaction \emph{means}, independent of which
screen happens to render it today, or which AI agent happens to emit it
tomorrow.

This is not a tooling gap; it is a categorical one. The SDV stack has rich
abstractions for signals (COVESA VSS), transports (uProtocol, Zenoh, SOME/IP),
data brokering (Kuksa), and middleware (AUTOSAR Classic and Adaptive). It has
rich vendor-specific frameworks for rendering (Qt Automotive, Kanzi, Android
Automotive). What it does not have is a layer \emph{between} them --- a
mediation boundary in which the meaning, attention demand, trust provenance,
and contextual fitness of an interaction are described independently of its
surface form. The interaction layer is where the value of the SDV becomes
perceptible, where regulatory exposure concentrates, and where the marginal
cost of change is highest --- yet it remains the least abstracted layer of the
stack.

Current HMI stacks are screen-first. A graphical layout binds a fixed widget
to a fixed input device; voice and haptic modalities are added as alternatives
after the fact. When a new surface arrives --- or an AI agent needs to emit an
interaction --- the logic must be rewritten. Engineering cost rises with each
new screen size or modality. User experience cost rises as the same intent
(\emph{acknowledge an alert}, \emph{increase volume}, \emph{navigate back})
acquires inconsistent behaviour across vehicles, generations, and contexts.
Security cost rises as AI agents and third-party applications gain the ability
to emit interactions that are indistinguishable, from the occupant's
standpoint, from safety-critical subsystems. These costs are symptoms of the
missing boundary --- not a tooling deficit, not a renderer deficit, but an
absent layer of architecture.

A vehicle in service for fifteen years will outlive its displays, its voice
models, its compute hardware, and quite possibly its operating system. What it
cannot afford to outlive is the \emph{meaning} of
\texttt{Alert.Collision.Warning}. We propose that meaning --- not pixels ---
is the right unit of design at the interaction layer, because meaning is what
must survive the turnover of every other component. The
\emph{Semantic Interaction Architecture} (SIA) is a vendor-neutral layer in
which selected in-vehicle interactions are described by their meaning,
attention demand, contextual fitness, and authority of origin, rather than by
buttons, screens, or widgets. Its central claim is uncomfortable but precise:
while the industry has spent two decades building the surface, the structure
underneath has remained implicit. SIA names it.

The first version of such a layer should be deliberately narrow. We scope the
proposal to three cores: \emph{intent/action abstraction} for high-value
commands and events, \emph{attention policy} for priority, interruptibility,
and driving context, and \emph{trust provenance} for determining which actors
may emit which interaction types with which authority.

%% ---------------------------------------------------------------
\section{Related Work}
%% ---------------------------------------------------------------

\paragraph{SDV data and service layers.}
COVESA VSS defines a vendor-neutral catalogue of vehicle signals~\cite{covesa}.
Eclipse Kuksa, uProtocol, and Chariott abstract data brokering, messaging, and
service discovery~\cite{eclipsesdv}. None of these projects model what an
interaction \emph{means to the occupant}; they model how services and signals
communicate.

\paragraph{HMI frameworks.}
Android Automotive's Car App Library is the closest production-grade analogue
to a semantic interaction model: applications declare templates and a host
renders them with built-in distraction optimisation~\cite{google-aaos}. Its
scope is bounded to supported app categories within the host-controlled
Android Automotive OS ecosystem.
Qt Automotive and Kanzi offer renderer-side abstractions without cross-renderer
interaction semantics.

\paragraph{Multimodal interaction.}
The W3C MMI Recommendation and EMMA annotation format define a generic semantic
model for multimodal input~\cite{w3c-mmi,w3c-emma}. Automotive uptake has been
limited; the vocabulary remains a candidate substrate for the input side of
SIA.

\paragraph{Adjacent domains.}
Unreal's Enhanced Input and Unity's Input System abstract actions from devices
in game engines. ARIA performs an analogous role for web accessibility. These
precedents demonstrate tractability; they have not been adapted to automotive
constraints (ASIL, regulated distraction limits, multi-renderer safety
requirements, 15-year vehicle lifecycles).

SIA occupies the gap above SDV service abstractions and below concrete
renderers --- the missing connective tissue that none of the above projects
addresses.

%% ---------------------------------------------------------------
\section{Mediation Architecture}
%% ---------------------------------------------------------------

We propose a mediation architecture of three functional components and two
cross-cutting policy functions. Renderers and input devices are external to SIA
and interface with it through capability declarations. This is not a complete
HMI platform; it defines what crosses the semantic boundary and how policy is
applied, leaving renderer implementation and GUI framework behaviour to
existing stacks.

\begin{figure}[h]
  \centering
  % Placeholder: export figures/fig2-mediation-architecture.pdf from the
  % Mermaid source in 01_Semantic-Interaction-Architecture-sdv.md (Fig 3 in the main paper)
  % and uncomment: \includegraphics[width=\columnwidth]{figures/fig2-mediation-architecture}
  \fbox{\parbox{0.92\columnwidth}{\centering\vspace{2em}
    \textit{[Figure 1: Mediation architecture — SIA contains three functional
    components (Ontology, Translation, Runtime) and two cross-cutting policies
    (Trust, Context). Emitters (Agents/ADAS) and renderers are external.
    Runtime flow: Emitters → Trust Policy → Translation → Runtime →
    external Renderers ↔ Occupant. Renderer capabilities feed back into
    Translation; Context Policy modulates Translation and Runtime.
    Insert PDF export of Mermaid diagram here.]}
  \vspace{2em}}}
  \caption{Mediation architecture. SIA contains three functional components
    (Ontology, Translation, Runtime) and two cross-cutting policies (Trust,
    Context). Emitters and renderers are external; they interface with SIA
    through Trust Policy (entry) and capability/dispatch flows (exit).}
  \label{fig:architecture}
\end{figure}

\paragraph{Ontology Language and Schema Profile.}
A stable language defines the long-term vocabulary of interaction meaning,
inheritance, metadata contracts, and compatibility rules. The first
standardisation target is a small typed event/command schema profile for
high-value interactions --- tractable to implement without locking in a deep
class hierarchy.

\paragraph{Translation Layer.}
A bidirectional adaptor that maps semantic nodes to concrete modalities given a
capability set and context. Inputs: the node, declared renderer capabilities,
the active context vector, and user accessibility profile. Output: a candidate
modality set and rendering or input mapping decision.

\paragraph{Interaction Coordination Runtime.}
A coordination function for focus, in-flight task flows, acknowledgement
timers, and consistency across distributed renderers. It does not replace a GUI
framework; it coordinates semantic state that multiple renderers must share.

\paragraph{Trust Policy.}
Validates all incoming nodes against the schema and trust requirements declared
in the Ontology Language. Trust Policy is a gate: a node that fails
verification is rejected and logged; it never reaches the Translation Layer.
The Ontology Language supplies the authoritative schema; the emitter supplies
the attestation; Trust Policy compares them.

\paragraph{Context Policy.}
Supplies a continuously updated context vector that modulates both the
Translation Layer (modality selection and attention budget) and the Interaction
Coordination Runtime (conflict resolution, acknowledgement timeouts).

\paragraph{Renderers and input devices --- external to SIA.}
Concrete output and input surfaces --- HUD, cluster, IVI touchscreen, voice,
haptic, AR, steering wheel controls --- are vendor-specific implementations
that consume modality decisions from the Runtime and declare measurable
capabilities into the Translation Layer. They are not components of SIA;
keeping them external preserves SIA's vendor neutrality. All nodes reaching
a renderer have already been verified by Trust Policy and coordinated by the
Runtime.

%% ---------------------------------------------------------------
\section{Node Taxonomy and Metadata Contracts}
%% ---------------------------------------------------------------

A common failure mode in interaction schemas is conflating semantically
different node types into a single generic message model. SIA separates four
primary semantic primitives, each with its own metadata contract.

\begin{figure}[h]
  \centering
  \fbox{\parbox{0.92\columnwidth}{\centering\vspace{2em}
    \textit{[Figure 2: Node taxonomy — Interaction splits into Action, Event,
    State, Task; Event splits into Alert and Notification.
    Insert PDF export of Mermaid diagram here.]}
  \vspace{2em}}}
  \caption{Node taxonomy. Four primary types with distinct metadata contracts;
    Event splits into Alert and Notification. Subclasses may strengthen but not
    weaken contracts.}
  \label{fig:taxonomy}
\end{figure}

\paragraph{Action} --- occupant-initiated; discrete, sustained, or continuous.
Carries \texttt{attention\_metrics}, \texttt{temporal\_type},
\texttt{recommended\_modality}.

\paragraph{Event} --- system-initiated; splits into \textbf{Alert}
(safety-relevant, may require acknowledgement) and \textbf{Notification}
(informational, suppressible). Alerts carry \texttt{priority},
\texttt{requires\_ack}, \texttt{trust\_requirements}.

\paragraph{State} --- focus, mode, and context transitions consumed by the
coordination runtime. Carries \texttt{scope}, \texttt{target\_role},
\texttt{consistency\_class}.

\paragraph{Task} --- a composed multi-step flow over Actions and States with
resumption semantics. Carries \texttt{step\_count},
\texttt{interruptible\_at}, \texttt{resumable\_across\_contexts}.

Every node carries two categories of metadata fields: \emph{declarative}
fields defined in the ontology and stable across deployments (e.g.,
\texttt{attention\_metrics}, \texttt{trust\_requirements},
\texttt{fallback\_chain}), and \emph{runtime} fields filled by the emitter at
the moment of emission (e.g., attestation, instance payload).
Table~\ref{tab:contracts} summarises mandatory, optional, and not-applicable
fields for each type.

\begin{table}[h]
  \caption{Partial metadata contract by node type. Full contract in the
    companion position paper.}
  \label{tab:contracts}
  \begin{tabular}{lccccc}
    \toprule
    Field & Action & Alert & Notif. & State & Task \\
    \midrule
    \texttt{since\_version}      & \textbullet & \textbullet & \textbullet & \textbullet & \textbullet \\
    \texttt{attention\_metrics}  & \textbullet & \textbullet & \textbullet & ---         & \textbullet \\
    \texttt{priority}            & ---         & \textbullet & \textbullet & ---         & ---         \\
    \texttt{requires\_ack}       & ---         & \textbullet & $\circ$     & ---         & ---         \\
    \texttt{trust\_requirements} & $\circ$     & \textbullet & $\circ$     & $\circ$     & $\circ$     \\
    \texttt{suppression\_class}  & ---         & \textbullet & \textbullet & ---         & ---         \\
    \texttt{fallback\_chain}     & \textbullet & \textbullet & \textbullet & ---         & \textbullet \\
    \texttt{degradation\_policy} & \textbullet & \textbullet & \textbullet & ---         & \textbullet \\
    \bottomrule
    \multicolumn{6}{l}{\small \textbullet\ mandatory \quad $\circ$\ optional \quad ---\ not applicable}
  \end{tabular}
\end{table}

%% ---------------------------------------------------------------
\section{Trust and Attention Policy}
%% ---------------------------------------------------------------

\subsection{Trust}

Current SDV cybersecurity practice --- ISO/SAE 21434~\cite{iso21434}, UNECE R155~\cite{unece-r155} --- addresses
firmware integrity, ECU authentication, and transport security. These are
necessary but insufficient: they do not constrain what an authenticated actor
is \emph{authorised to say} at the interaction level. An attacker who cannot
compromise the brakes can still suppress a collision warning, inject a fake
alert, or elevate the priority of a benign notification to displace a critical
one.

SIA separates trust into two artefacts. \emph{Trust requirements} are declared
on the node in the ontology --- what consumers must verify before propagating
or rendering. \emph{Trust attestation} is attached to the instance by the
emitter --- cryptographic and provenance evidence that requirements are met.
The Trust Policy verifies attestation satisfies requirements before the node
propagates. Trust failures are \textbf{fail-closed}: the node is rejected and a
\texttt{SecurityEvent} is logged; it never reaches Translation or a renderer.
This is distinct from \texttt{degradation\_policy}, which governs only renderer
capability fallback (e.g., routing to voice when the HUD is unavailable)
\emph{after} a trust check has already passed.

An \texttt{actor\_class} taxonomy drives policy mechanically: \texttt{adas} and
\texttt{vsc} may emit safety-critical alerts; \texttt{agent\_local} and
\texttt{agent\_cloud} may not --- regardless of how the agent was prompted.
This is a structurally stronger property than service-level authentication: it
constrains what \emph{kinds of things} an authenticated actor may say, not
merely whether it may speak.

\subsection{Attention}

Contemporary HMI guidelines (NHTSA~\cite{nhtsa2013}, ISO 15005~\cite{iso15005}, JAMA~\cite{jama2004}) prescribe
measurable thresholds --- total eyes-off-road time (TEORT), single glance
duration, task step count. Most software-side frameworks operate on qualitative
tags (\texttt{distraction-optimised: true/false}) not directly auditable
against those thresholds.

SIA attaches predicted attention-demand proxies to every interaction-bearing
node:

\begin{lstlisting}[language={}]
attention_metrics:
  glance_time_estimated_ms: 1500
  mean_single_glance_ms: 400
  task_steps: 3
  voice_alt_available: true
  cognitive_load: moderate
\end{lstlisting}

The Translation Layer composes node metrics with a context modifier from
Context Policy:
\[
  \text{effective\_cost} = \text{base\_metric} \times \text{context\_modifier}
\]
Budgets are configured \emph{per node class and per context}, not globally.
For illustration: $\le$1500~ms TEORT for \texttt{Alert.*} (safety-critical,
must surface immediately) and $\le$2000~ms TEORT for general \texttt{Action.*}
under manual highway driving. Concrete budgets are a deployment-level
configuration; SIA defines only the contract that makes them mechanically
enforceable.

%% ---------------------------------------------------------------
\section{Worked Example: Alert.Collision.Warning}
%% ---------------------------------------------------------------

We trace a single safety-critical alert end-to-end to make the contracts
concrete.

\paragraph{Ontology declaration (abbreviated):}

\begin{lstlisting}[language={}]
node: Interaction.Event.Alert.Collision.Warning
priority: 95
requires_ack: true
ack_authority: driver_only
trust_requirements:
  permitted_actor_classes: [adas, vsc]
  max_age_ms: 200
  replay_protection: required
attention_metrics:
  glance_time_estimated_ms: 600
suppression_class: safety_critical
fallback_chain: [hud, cluster, voice, haptic]
\end{lstlisting}

\paragraph{Trust verification:}

\begin{figure}[h]
  \centering
  \fbox{\parbox{0.92\columnwidth}{\centering\vspace{2em}
    \textit{[Figure 3: Sequence diagram — ADAS emits instance + attestation to
    Trust Policy; on failure: reject + log SecurityEvent; on success: verified
    node propagates to Translation, Runtime, Renderers; ack closes loop.
    Insert PDF export of Mermaid sequenceDiagram here.]}
  \vspace{2em}}}
  \caption{Sequence flow for Alert.Collision.Warning. Trust Policy is a
    chokepoint before Translation. Renderer dispatch is multicast;
    acknowledgement is tracked by Runtime.}
  \label{fig:sequence}
\end{figure}

\paragraph{Context-dependent translation.}
Under manual highway driving (\texttt{sae\_level: 1},
\texttt{vehicle\_state: moving}, \texttt{traffic\_density: dense}), the
Translation Layer selects HUD as primary with concurrent cluster and haptic,
and rejects the IVI touchscreen (off-axis; exceeds attention budget under
dense-traffic context modifier). Under L4 autonomous driving
(\texttt{sae\_level: 4}, \texttt{driver\_state: not\_monitoring}), HUD is
de-prioritised in favour of cluster and full-sentence voice prompt; Context
Policy scales the effective \texttt{ack\_timeout\_ms} from a base of 3000~ms
to 6000~ms via its published modifier rule (the node's declarative value is
unchanged).

\paragraph{Adversarial rejection.}
A third-party application addressing \texttt{Alert.Collision.Warning} is
rejected as an unauthorised emission
(\texttt{third\_party\_app} $\notin$ \texttt{permitted\_actor\_classes}), even
when its own app-store signature is valid. A class-spoofing attempt --- the
same app falsely attesting \texttt{actor\_class: adas} --- fails at signature
verification because the app lacks the ADAS signing key. A replay of a
1100~ms-old instance is rejected because age exceeds
\texttt{max\_age\_ms: 200}. A local LLM agent is rejected on the same class
basis. Priority injection --- an adversary emitting a notification with
\texttt{priority: 99} --- is defeated because priority is a property of the
ontology declaration, not the instance.

%% ---------------------------------------------------------------
\section{Implications for AutomotiveUI Research}
%% ---------------------------------------------------------------

SIA's value depends on empirical foundations that the AutomotiveUI community
is uniquely positioned to provide. We identify five concrete research handles,
each addressable with methods this community already practices.

\paragraph{Attention as a first-class research variable.}
The \texttt{attention\_metrics} contract makes declared attention cost
comparable across HMI designs, OEMs, and evaluation studies --- without
requiring every comparison to be grounded in a new user study. Researchers
could evaluate whether declared proxies predict empirical TEORT~\cite{ebel2021},
and how context modifiers should be calibrated. This creates a feedback loop
between user studies and the ontology itself.

\paragraph{Cross-renderer semantic equivalence.}
By expressing the same interaction node across modalities (visual, voice,
haptic), SIA enables controlled studies of cross-modal equivalence --- does
\texttt{Alert.Collision.Warning} delivered via voice versus HUD produce
equivalent driver response? Currently such comparisons are confounded by
differing implementations; a shared semantic node removes that confound.

\paragraph{Trust and agency in AI-augmented vehicles.}
The \texttt{actor\_class} taxonomy provides a vocabulary for the growing
AutomotiveUI literature on in-vehicle agents~\cite{demir2024,gomaa2022}.
Empirical work on driver trust calibration, appropriate reliance, and agent
transparency can be grounded in the structural distinction between
\texttt{agent\_local}, \texttt{agent\_cloud}, \texttt{adas}, and
\texttt{human\_direct} --- classes that carry different implied authority and
verifiability.

\paragraph{Context-adaptive interaction design.}
The multi-axis context vector (\texttt{sae\_level}, \texttt{driver\_state},
\texttt{road\_type}, \texttt{vehicle\_state}, \texttt{market\_jurisdiction})
formalises the context space that AutomotiveUI researchers use informally. Experiments can be designed
against explicit context predicates rather than bespoke scenario definitions,
improving replicability.

\paragraph{Standardisation input.}
AutomotiveUI researchers regularly engage with industry and standards bodies.
Empirical validation of the attention metric composition formula, the priority
scale, and the \texttt{ack\_kind} distinction (certified versus inferred
acknowledgement) are directly actionable contributions to a future
standardisation process.

None of these threads requires the architecture be standardised first; on the
contrary, the standardisation effort itself depends on this work happening
early. We invite collaboration on any of them --- from individual empirical
studies to multi-site replications --- and treat AutomotiveUI as the principal
venue where SIA's empirical foundations should be developed.

%% ---------------------------------------------------------------
\section{Open Questions}
%% ---------------------------------------------------------------

This work is scoped to a position paper. We surface the following as the most
urgent open questions, where community contribution would directly shape the
standard's first version.

\begin{enumerate}
  \item \emph{Schema formalism.} JSON Schema, OWL/SHACL, or a VSS-style custom
    format generating to JSON Schema. The choice has long-term tooling and
    adoption consequences.

  \item \emph{Empirical validation of attention metric composition.} The
    $\text{effective\_cost} = \text{base} \times \text{context\_modifier}$
    formula requires user study evidence. The composition law is likely not
    scalar.

  \item \emph{Certified versus inferred acknowledgement.}
    \texttt{ack\_kind: gaze} mixes deterministic input with probabilistic
    inference; the ontology may need to distinguish certified from inferred
    acknowledgement, with different safety implications.

  \item \emph{Multi-hop trust provenance.} In agentic deployments, the
    provenance chain may have multiple hops (\texttt{sensor} $\to$
    \texttt{adas} $\to$ \texttt{trust\_verifier} $\to$ \texttt{runtime});
    chain-level trust composition is not yet specified.

  \item \emph{Reference implementation and tooling roadmap.} Architecture
    adoption is gated not only on specification quality but on the availability
    of a supporting toolchain. Prior art from W3C MMI and related efforts
    suggests that well-specified interfaces without reference tooling
    consistently fail to achieve adoption. We identify three tooling layers,
    each a prerequisite for the next.

    \begin{description}
      \item[Core layer.]
        A machine-readable schema (JSON Schema or Protobuf) for Action, Event,
        State, and Task nodes with their full metadata contracts; a validation
        engine that answers \emph{``may this node be rendered in this
        context?''} given the active context vector; and a declarative policy
        engine (OPA/Rego-style) encoding rules such as
        \emph{``\texttt{agent\_local} actors may not emit
        \texttt{safety\_critical} nodes''} and \emph{``a
        \texttt{fallback\_chain} with no \texttt{voice} entry is invalid during
        L4 autonomy.''}

      \item[Integration layer.]
        Thin adaptors to COVESA VSS, Eclipse Kuksa, and
        uProtocol~\cite{eclipsesdv} so that SIA is a plug-in layer over
        existing SDV transports rather than a replacement. A renderer capability
        interface --- analogous to Android's capability detection or React
        Native's platform abstraction --- through which concrete surfaces
        declare what they can render and at what attention cost. An agent
        interface that accepts semantic node payloads from LLM-based agents and
        routes them through trust verification before any UI surface is reached.

      \item[Developer tooling.]
        A context simulator that exercises the Translation Layer against
        scripted driving contexts and failure modes; a structured debugger that
        makes suppression decisions transparent (\emph{``collision warning
        suppressed --- reason: \texttt{actor\_class agent\_local}
        $\notin$ \texttt{permitted\_actor\_classes}''});
        and a declarative test harness expressing automotive certification
        scenarios as \texttt{given}/\texttt{when}/\texttt{then} predicates
        against the schema and policy engine. The test harness is likely the
        most direct contribution to ASIL-relevant certification workflows and
        the tool most likely to determine whether practitioners adopt SIA or
        not.
    \end{description}
\end{enumerate}

Each of these is approachable through methods already standard in this
community: simulator studies for attention calibration (\#2), controlled
experiments for the certified-vs-inferred acknowledgement distinction (\#3),
mixed-methods work on multi-hop trust scenarios (\#4), and
reference-implementation contributions for the tooling roadmap (\#5). Short
empirical studies and full specification contributions are equally welcome ---
and equally needed to shape the standard before it freezes.

%% ---------------------------------------------------------------
%% References
%% ---------------------------------------------------------------
\bibliographystyle{ACM-Reference-Format}
\bibliography{04_WiP-references}

\end{document}
